\documentclass[11pt,a4paper]{article}

\usepackage{amsmath,amssymb,amsfonts}
\usepackage{graphicx}
\usepackage{hyperref}
\usepackage{booktabs}
\usepackage{amsthm}
\usepackage[margin=2.5cm]{geometry}

\newcommand{\Spec}{\mathrm{Spec}}
\newcommand{\Zp}{\mathbb{Z}[1/p]}

\begin{document}

\title{Physical Implementation of Prime Channel Muting:\\
  Ten Experimental Platforms for Arithmetic Vacuum Engineering}

\author{Wright Brothers\\{\small Independent Research}}
\date{April 2026}
\maketitle

\begin{abstract}
The companion paper~\cite{paper_A} established that suppressing
vacuum modes at multiples of a prime $p$ (``prime channel muting'')
flips the sign of zeta-regularized vacuum energy, and proposed a
single experimental platform (superconducting circuits with SQUID
notch filters). The present paper systematically surveys
\emph{all} viable physical platforms for implementing prime
channel muting, organized by the three mechanisms identified in
the Bost-Connes framework: (A)~spectral filtering,
(B)~arithmetic boundary conditions, and (C)~phase transition
sector selection. We identify ten distinct implementation
candidates, assess each on five criteria (technology readiness,
number of primes addressable, directness of vacuum energy
measurement, cost, and timeline), and propose three
experiments executable with minimal resources: an acoustic
resonator proof-of-concept (\$100, 1~day), a quantum computer
simulation of the Bost-Connes Hamiltonian (\$0, 1~week on
cloud quantum hardware), and an optical frequency comb
experiment with Fabry-P\'erot filtering (\$10K, 6~months).
These complement the SQUID experiment of~\cite{paper_A} by
providing independent tests of the arithmetic structure of
vacuum energy across different physical domains.
\end{abstract}

% ============================================================================
\section{Introduction}
\label{sec:intro}
% ============================================================================

In~\cite{paper_A}, we proved that muting any single prime $p$
flips the sign of zeta-regularized vacuum energy:
$\zeta_{\neg p}(-3) = (1-p^3)/120 < 0$.
Three physical mechanisms were identified for implementing the
projection $P_{\neg p} = \sum_{p \nmid n} |n\rangle\langle n|$:
\begin{itemize}
\item[(A)] Spectral filtering: absorb modes at frequencies
  that are multiples of $p$.
\item[(B)] Arithmetic boundary conditions: impose boundary
  conditions that permit only modes $n$ with $\gcd(n,p)=1$.
\item[(C)] Phase transition sector selection: select KMS states
  of the Bost-Connes system invariant under $\mathbb{Z}_p^*$.
\end{itemize}

Paper~\cite{paper_A} developed mechanism~(A) into a detailed
SQUID-based experiment. The present paper surveys all mechanisms
systematically and identifies nine additional platforms, including
three that can be executed immediately with minimal resources.

% ============================================================================
\section{Mechanism A: Spectral Filtering}
\label{sec:mech_a}
% ============================================================================

Mechanism~A suppresses modes at multiples of $p$ by absorbing
the corresponding frequency components. We distinguish two
sub-cases depending on the spectrum of the physical system.

\subsection{A1: Logarithmic spectrum systems}

In the Bost-Connes system, $E_n = \log(n)$. A single notch
filter at $\omega = \log(p)$ removes all $p$-divisible states,
because any $n$ with $p \mid n$ has $\log(n) \ni \log(p)$ in
its prime factorization contribution.

\textbf{Platform 1: Charged wire with logarithmic potential.}
A 2D charged wire creates potential $V(r) = V_0 \log(r/a)$.
In the WKB approximation, the bound-state energies scale as
$E_n \sim \log(n)$. A particle confined in this potential with
a frequency-selective absorber at $\omega = \log(p)$ would
realize the BC Hamiltonian with prime muting.

Technology readiness: low (requires custom trap geometry).
Timeline: 3--5 years.

\subsection{A2: Uniform spectrum systems}

Standard cavities have $\omega_n = n\omega_0$ (uniform spacing).
To mute prime $p$, one must suppress modes at $p\omega_0$,
$2p\omega_0$, $3p\omega_0$, etc.\ --- a \emph{periodic comb filter}
with period $p\omega_0$.

\textbf{Platform 2: Superconducting circuit with SQUID filters}
(detailed in~\cite{paper_A}). Coplanar waveguide resonator
($f_0 \approx 6$\,GHz) with SQUID-based notch filters tuned to
each even harmonic. Technology readiness: high.
Cost: \$50--100K. Timeline: 12 months.

\textbf{Platform 3: Optical frequency comb with Fabry-P\'erot.}
A mode-locked laser generates a frequency comb
$f_n = f_{\mathrm{CEO}} + n f_{\mathrm{rep}}$ with
$f_{\mathrm{rep}} \sim 100$\,MHz.
A Fabry-P\'erot cavity with free spectral range
$\mathrm{FSR} = p \cdot f_{\mathrm{rep}}$ transmits only modes
at multiples of $p$ and reflects all others.

The \emph{reflected} beam contains exactly the modes coprime to $p$:
$\{n : p \nmid n\}$. This is the $p$-muted spectrum.

Measurement: the quantum noise (shot noise) of the reflected beam
encodes the vacuum fluctuations of the filtered mode set.
Comparing noise power between filtered and unfiltered beams
tests whether vacuum fluctuation energy depends on the arithmetic
of mode selection.

Technology readiness: high (frequency combs are Nobel Prize
technology~\cite{Hansch2005}; Fabry-P\'erot cavities are standard).
Cost: \$10--50K. Timeline: 6 months.

\textbf{Platform 4: Acoustic resonator with absorbers.}
A tube of length $L$ supports standing waves at $f_n = nv/(2L)$.
Inserting an absorbing material at position $x = L/(2p)$
preferentially damps modes whose spatial profile has an antinode
at that position.

For $p = 2$: absorber at $L/4$ damps odd modes (antinode)
while even modes have a node there and survive.
This gives the \emph{complement} of $p$-muting (even modes survive
instead of being removed), but the principle is identical and
the complementary measurement is equally informative.

Technology readiness: trivial (acrylic tube, speaker, microphone).
Cost: \$100. Timeline: 1 day. This is the minimum-cost
proof-of-concept for arithmetic mode selection.

\textbf{Platform 5: Piezoelectric resonator array.}
Piezoelectric crystals (quartz, PZT) have sharp mechanical
resonances. An array of $N$ crystals with resonant frequencies
$f_n = n f_0$ and selective damping at multiples of $p$ provides
a solid-state platform for acoustic prime muting.

Technology readiness: moderate. Cost: \$5K. Timeline: 3 months.

% ============================================================================
\section{Mechanism B: Arithmetic Boundary Conditions}
\label{sec:mech_b}
% ============================================================================

Mechanism~B imposes boundary conditions that geometrically
select modes coprime to $p$, analogous to how conducting plates
select modes fitting between them in the Casimir effect.

\textbf{Platform 6: Semiconductor superlattice.}
A GaAs/AlGaAs superlattice with period $d = p \cdot a_0$
(where $a_0$ is the fundamental lattice constant) opens
minibands with gaps at mode numbers that are multiples of $p$.
Electrons in the allowed minibands occupy only modes
with $\gcd(n, p) = 1$.

For $p = 2$: $a_0 = 5$\,nm, $d = 10$\,nm, gap energy $\sim 0.1$\,eV.
Even-numbered minibands are gapped; odd ones propagate.

Technology readiness: high (MBE growth is mature technology).
Cost: \$50K. Timeline: 1 year.

\textbf{Platform 7: Electromagnetic metamaterial.}
A 3D metamaterial with split-ring resonators at period $p \cdot a$
opens electromagnetic band gaps at multiples of $p$.
Three independent lattice directions can address three independent
primes: period $2a$ along $x$, $3a$ along $y$, $5a$ along $z$
simultaneously mutes $p = 2, 3, 5$.

Technology readiness: moderate. Cost: \$200K. Timeline: 2 years.

\textbf{Platform 8: Euler product crystal.}
The most ambitious version of Mechanism~B: stack multiple
Bragg gratings, each with a different period $p_i \cdot a$
for primes $p_1, p_2, p_3, \ldots$

Layer 1 (period $2a$): gaps at multiples of 2.
Layer 2 (period $3a$): gaps at multiples of 3.
Layer 3 (period $5a$): gaps at multiples of 5.

The combination is a physical \emph{sieve of Eratosthenes}:
only prime-numbered modes (plus $n = 1$) survive all layers.
The effective zeta function becomes
$\zeta(s) \prod_{i} (1 - p_i^{-s})$, removing multiple
Euler factors simultaneously.

Numerical estimate: with the first three primes ($2, 3, 5$),
77\% of modes below $n = 100$ are removed.

Technology readiness: low (multi-layer interference is challenging).
Cost: \$500K. Timeline: 3--5 years.

% ============================================================================
\section{Mechanism C: Phase Transition Sector Selection}
\label{sec:mech_c}
% ============================================================================

Mechanism~C implements prime muting by manipulating the
Bost-Connes phase transition directly, selecting KMS states
invariant under $\mathbb{Z}_p^*$.

\textbf{Platform 9: Quantum computer simulation.}
The BC Hamiltonian $H|n\rangle = \log(n)|n\rangle$ on
$\ell^2(\{1, \ldots, 2^N\})$ can be implemented on $N$ qubits.
The prime-muting perturbation
$H' = H + \Delta \sum_{p|n} |n\rangle\langle n|$
gaps out $p$-divisible states in the limit $\Delta \to \infty$.

Protocol:
(1) Prepare the ground state of $H$ via VQE or adiabatic preparation.
(2) Measure $\langle H \rangle$.
(3) Add perturbation $H'$ with increasing $\Delta$.
(4) Measure $\langle H' \rangle$ as a function of $\Delta$.
(5) Verify that $\langle H' \rangle \to \zeta_{\neg p}(-1)$
in the $\Delta \to \infty$ limit.

This can be run on IBM Quantum or Google Cirq cloud platforms
with 5--10 qubits, covering the first 32--1024 modes.

Technology readiness: high (quantum simulation is standard).
Cost: \$0 (cloud free tier). Timeline: 1 week.

\textbf{Platform 10: Cold atom optical lattice.}
Load ultracold atoms into a 1D optical lattice with
site-dependent potential $V_n = V_0 \log(n)$
(approximated by a shaped laser beam).
A secondary lattice of period $p$ with depth $\Delta$
shifts the energy of $p$-divisible sites.

For $p = 2$: the secondary lattice has wavelength
$\lambda_2 = 2 \times \lambda_1 = 1064$\,nm (Nd:YAG fundamental).
For $p = 3$: $\lambda_3 = 3 \times \lambda_1 = 1596$\,nm.
For $p = 5$: $\lambda_5 = 5 \times \lambda_1 = 2660$\,nm.
All are standard laser wavelengths.

The cold-atom platform can address thousands of lattice sites
and multiple primes simultaneously, making it the most scalable
analog quantum simulator for the BC system.

A closely related platform uses Bose-Einstein condensates (BECs)
with spatially modulated speed of sound. Since BEC phonons
realize ``analog gravity''~\cite{Unruh1981}, a BEC with
arithmetic modulation would implement
\emph{analog gravity in an arithmetically modified vacuum} ---
the most direct bridge from our algebraic mechanism to
spacetime physics.

Technology readiness: moderate (optical lattice experiments
are mature; the logarithmic potential is non-standard).
Cost: \$500K (existing lab) to \$2M (new setup).
Timeline: 2--3 years.

% ============================================================================
\section{Comparative Assessment}
\label{sec:comparison}
% ============================================================================

Table~\ref{tab:comparison} summarizes all ten platforms.

\begin{table*}[t]
\centering
\caption{Ten implementation platforms for prime channel muting.
  TRL = technology readiness (1--5 scale).
  $N_p$ = number of primes simultaneously addressable.
  Measurement: D = direct vacuum energy, I = indirect (spectrum).}
\label{tab:comparison}
\small
\begin{tabular}{@{}rlcccccc@{}}
\toprule
\# & Platform & Mechanism & TRL & $N_p$ & Meas.\ & Cost & Timeline \\
\midrule
1 & Logarithmic potential trap & A1 & 2 & 1--2 & I & \$100K & 3--5 yr \\
2 & SC circuit + SQUID~\cite{paper_A} & A2 & 4 & 2--5 & D & \$100K & 1 yr \\
3 & Frequency comb + Fabry-P\'erot & A2 & 4 & 1--3 & I & \$10--50K & 6 mo \\
4 & Acoustic resonator & A2 & 5 & 1--2 & I & \$100 & 1 day \\
5 & Piezoelectric array & A2 & 4 & 1--2 & I & \$5K & 3 mo \\
6 & Semiconductor superlattice & B & 4 & 1--2 & I & \$50K & 1 yr \\
7 & EM metamaterial & B & 3 & 1--3 & I & \$200K & 2 yr \\
8 & Euler product crystal & B & 2 & 3--5 & I & \$500K & 3--5 yr \\
9 & Quantum computer (cloud) & C & 5 & many & I & \$0 & 1 wk \\
10 & Cold atom optical lattice & C & 3 & 3--5 & D & \$500K & 2--3 yr \\
\bottomrule
\end{tabular}
\end{table*}

% ============================================================================
\section{Three Immediately Executable Experiments}
\label{sec:immediate}
% ============================================================================

We highlight three experiments that can begin immediately,
complementing the SQUID experiment of~\cite{paper_A}.

\subsection{Experiment $\alpha$: Acoustic resonator (\$100, 1 day)}

\textbf{Materials}: acrylic tube ($L = 1$\,m), loudspeaker, microphone, PC with FFT software.

\textbf{Protocol}:
(1) Drive the tube with white noise.
(2) Record the resonance spectrum via FFT: peaks at
$f_n = nv/(2L) \approx n \times 172$\,Hz.
(3) Insert absorbing foam at $L/2$: odd modes (antinode at $L/2$)
are damped; even modes (node at $L/2$) survive.
(4) Insert foam at $L/3$: modes not divisible by 3 that have
an antinode there are further damped.
(5) Compare spectral peak heights before and after insertion.

\textbf{What it tests}: That physical absorbers at arithmetic
positions selectively remove modes based on divisibility.
This is a classical (non-quantum) demonstration of
arithmetic mode selection.

\textbf{Limitation}: Does not probe vacuum energy
(classical acoustics). But establishes the principle that
physical geometry implements arithmetic filtering.

\subsection{Experiment $\beta$: Quantum simulation (\$0, 1 week)}

\textbf{Platform}: IBM Quantum (free cloud access) or Google Cirq.

\textbf{Protocol}:
(1) Encode the BC Hamiltonian $H|n\rangle = \log(n)|n\rangle$
on 5--10 qubits using diagonal gates.
(2) Prepare the thermal state at inverse temperature $\beta$
via variational quantum eigensolver (VQE).
(3) Measure $\langle H \rangle = -\partial_\beta \log Z(\beta)$.
(4) Add perturbation $\Delta \sum_{p|n} |n\rangle\langle n|$
for $p = 2$ with increasing $\Delta$.
(5) Measure $\langle H \rangle$ as function of $\Delta$.
(6) Verify convergence to $\zeta_{\neg 2}(-1) = +1/12$ as
$\Delta \to \infty$ (sign flip from $\zeta(-1) = -1/12$).

\textbf{What it tests}: That the BC system with prime muting
reproduces the predicted partition function modification.
This is a \emph{digital} quantum simulation, not an analog
experiment, but validates the algebraic mechanism on
actual quantum hardware.

\textbf{Limitation}: Simulates the BC system, not the
physical vacuum. Does not prove that physical vacuum
has BC structure.

\subsection{Experiment $\gamma$: Frequency comb (\$10--50K, 6 months)}

\textbf{Materials}: Mode-locked laser ($f_{\mathrm{rep}} \sim 100$\,MHz),
Fabry-P\'erot cavity ($\mathrm{FSR} = 2 f_{\mathrm{rep}}$),
balanced homodyne detector.

\textbf{Protocol}:
(1) Generate frequency comb: $f_n = f_{\mathrm{CEO}} + n f_{\mathrm{rep}}$.
(2) Pass through Fabry-P\'erot with $\mathrm{FSR} = 2 f_{\mathrm{rep}}$:
transmitted beam contains only even modes ($p = 2$ multiples);
reflected beam contains only odd modes ($p = 2$ muted).
(3) Measure quantum noise power of reflected beam
(shot noise limit = vacuum fluctuation energy).
(4) Compare to unfiltered beam noise.
(5) Repeat with $\mathrm{FSR} = 3 f_{\mathrm{rep}}$ (mute $p = 3$)
and $\mathrm{FSR} = 5 f_{\mathrm{rep}}$ (mute $p = 5$).
(6) Test multiplicative structure:
noise ratio $R(p_1, p_2) = R(p_1) \times R(p_2)$?

\textbf{What it tests}: Whether vacuum fluctuation noise of
a filtered mode set depends on the arithmetic of the
filter in a manner consistent with the Euler product.

\textbf{Limitation}: Measures vacuum \emph{fluctuations},
not vacuum \emph{energy} directly. The connection between
noise power and zeta-regularized energy requires the
quantum optics formalism relating shot noise to zero-point
energy~\cite{Caves1981}.

% ============================================================================
\section{Scaling Considerations}
\label{sec:scaling}
% ============================================================================

In~\cite{paper_B}, we conjectured that the weak energy condition
is a discrete topological invariant (regulator orientation)
rather than a continuous energy threshold. If correct, scaling
is not needed: a single localization $\mathbb{Z} \to \Zp$
flips the WEC regardless of energy scale.

Independent of this conjecture, we note that each platform
has a natural scaling path:
\begin{itemize}
\item Platform 2 (SQUID): arrays of $N$ synchronized resonators,
  $E \propto N^2$ if vacuum energy modifications are coherent.
\item Platform 6 (superlattice): bulk material scales with volume.
\item Platform 8 (Euler product crystal): 3D structure, energy
  density scales with the number of Euler factors removed.
\item Platform 10 (optical lattice/BEC): $\sim 10^6$ sites
  in current experiments; phonon-mediated interactions provide
  natural long-range coherence.
\end{itemize}

The BEC platform (variant of Platform~10) is particularly
promising for scaling because: (a) BEC phonons are an
established analog gravity platform~\cite{Unruh1981};
(b) Feshbach resonances allow precise tuning of the speed
of sound, enabling spatial modulation at arithmetic periods;
(c) the number of modes in a BEC ($\sim 10^6$) far exceeds
what is achievable in microwave circuits ($\sim 20$).

% ============================================================================
\section{Discussion}
\label{sec:discussion}
% ============================================================================

The ten platforms span four orders of magnitude in cost
(\$0 to \$2M) and five orders in timeline (1 day to 5 years).
They address all three mechanisms (A, B, C) and both spectral
regimes (microwave and optical).

We recommend a multi-pronged experimental strategy:
\begin{enumerate}
\item \textbf{Immediate} (month 1):
  Execute Experiments $\alpha$ (acoustic, principle verification)
  and $\beta$ (quantum simulation, BC system verification)
  simultaneously. Cost: \$100 + \$0 = \$100.
\item \textbf{Short-term} (months 1--6):
  Pursue Experiment $\gamma$ (frequency comb, vacuum fluctuation test).
  Cost: \$10--50K.
\item \textbf{Medium-term} (months 6--18):
  Execute the SQUID experiment of~\cite{paper_A}
  (Platform 2, direct vacuum energy measurement).
  Cost: \$50--100K.
\item \textbf{Long-term} (years 2--5):
  Develop semiconductor superlattice (Platform 6) or
  cold atom (Platform 10) for multi-prime muting and scaling studies.
\end{enumerate}

If Experiments $\alpha$ and $\beta$ succeed (confirming the
principle and the BC algebraic structure), and Experiment $\gamma$
shows arithmetic dependence of vacuum noise, the case for the
more expensive Platforms 2, 6, and 10 becomes compelling.

% ============================================================================
\section{Conclusion}
\label{sec:conclusion}
% ============================================================================

We have identified ten physical platforms for implementing
prime channel muting, spanning acoustic, electromagnetic,
superconducting, semiconductor, metamaterial, quantum computing,
and cold-atom technologies. Three experiments can begin
immediately at minimal cost (\$100, \$0, and \$10K respectively),
providing independent tests of the arithmetic vacuum energy
hypothesis across classical acoustics, digital quantum simulation,
and quantum optics.

The diversity of platforms is a strength of the arithmetic
vacuum engineering program: if the Euler product structure
of vacuum energy is a physical law (not just a regularization
artifact), it should manifest across \emph{all} physical systems
with quantized mode spectra, regardless of the specific
physical substrate. Multi-platform confirmation would provide
the strongest possible evidence.

\begin{thebibliography}{20}

\bibitem{paper_A}
Wright Brothers,
``Arithmetic vacuum engineering: Prime channel muting via
Bost-Connes phase transitions,''
companion paper (2026).

\bibitem{paper_B}
Wright Brothers,
``Sheaf-theoretic reformulation of vacuum energy conditions
on $\Spec(\mathbb{Z})$,''
companion paper (2026).

\bibitem{Hansch2005}
T.~W.~H\"ansch,
``Nobel Lecture: Passion for precision,''
Rev. Mod. Phys. \textbf{78}, 1297 (2006).

\bibitem{Unruh1981}
W.~G.~Unruh,
``Experimental black-hole evaporation?''
Phys. Rev. Lett. \textbf{46}, 1351 (1981).

\bibitem{Caves1981}
C.~M.~Caves,
``Quantum-mechanical noise in an interferometer,''
Phys. Rev. D \textbf{23}, 1693 (1981).

\bibitem{Wilson2011}
C.~M.~Wilson et al.,
Nature \textbf{479}, 376 (2011).

\end{thebibliography}

\end{document}
